\documentclass[11pt]{article}
\usepackage[utf8]{inputenc}
\usepackage[T1]{fontenc}     % Fix weird character
\usepackage{geometry}
\usepackage{amsmath}
\usepackage{amssymb}
\usepackage{gensymb}
\usepackage{spalign}
\usepackage{xfrac}
\usepackage{parskip}
\usepackage{float}  % figure[H]
\usepackage[english]{babel}
\usepackage[style=ieee,backend=biber,urldate=iso,date=iso]{biblatex}
\usepackage[breaklinks=true,bookmarks=true,hidelinks]{hyperref}
\usepackage{tikz}
\usepackage{pgfplots}
\usepackage{circuitikz}
\usepackage{subcaption}
\usepackage[yyyymmdd]{datetime}
\usepackage{url}
\usetikzlibrary {arrows.meta}
\pgfplotsset{compat=newest,compat/show suggested version=false}

% Custom shape
\makeatletter
\pgfdeclareshape{dflowpass}{
    \inheritsavedanchors[from=rectangle]
    \inheritanchorborder[from=rectangle]
    \inheritanchor[from=rectangle]{center}
    \anchor{left}{\pgfpoint{-0.5cm}{0cm}}
    \anchor{right}{\pgfpoint{0.5cm}{0cm}}
    % Add more custom anchors if needed
    \backgroundpath{
        % Draw the switch shape (rectangle with diagonal lines)
        \pgfsetlinewidth{0.85pt}
        \pgfpathrectanglecorners{\pgfpoint{-0.5cm}{-0.5cm}}{\pgfpoint{0.5cm}{0.5cm}}
        \pgfpathmoveto{\pgfpoint{-0.35cm}{0.125cm}}
        \pgfpathlineto{\pgfpoint{0.15cm}{0.125cm}}
        \pgfpathlineto{\pgfpoint{0.35cm}{-0.2cm}}
        \pgfusepath{draw}
    }
}
\pgfdeclareshape{dfhighpass}{
    \inheritsavedanchors[from=rectangle]
    \inheritanchorborder[from=rectangle]
    \inheritanchor[from=rectangle]{center}
    \anchor{left}{\pgfpoint{-0.5cm}{0cm}}
    \anchor{right}{\pgfpoint{0.5cm}{0cm}}
    \backgroundpath{
        % Draw the switch shape (rectangle with diagonal lines)
        \pgfsetlinewidth{1pt}
        \pgfpathrectanglecorners{\pgfpoint{-0.5cm}{-0.5cm}}{\pgfpoint{0.5cm}{0.5cm}}
        \pgfpathmoveto{\pgfpoint{0.35cm}{0.125cm}}
        \pgfpathlineto{\pgfpoint{-0.15cm}{0.125cm}}
        \pgfpathlineto{\pgfpoint{-0.35cm}{-0.2cm}}
        \pgfusepath{draw}
    }
}
\pgfdeclareshape{dfbandpass}{
    \inheritsavedanchors[from=rectangle]
    \inheritanchorborder[from=rectangle]
    \inheritanchor[from=rectangle]{center}
    \anchor{left}{\pgfpoint{-0.5cm}{0cm}}
    \anchor{right}{\pgfpoint{0.5cm}{0cm}}
    \backgroundpath{
        % Draw the switch shape (rectangle with diagonal lines)
        \pgfsetlinewidth{1pt}
        \pgfpathrectanglecorners{\pgfpoint{-0.5cm}{-0.5cm}}{\pgfpoint{0.5cm}{0.5cm}}
        \pgfpathmoveto{\pgfpoint{-0.35cm}{-0.2cm}}
        \pgfpathlineto{\pgfpoint{-0.20cm}{0.125cm}}
        \pgfpathlineto{\pgfpoint{0.15cm}{0.125cm}}
        \pgfpathlineto{\pgfpoint{0.35cm}{-0.2cm}}
        \pgfusepath{draw}
    }
}
\pgfdeclareshape{quantizer}{
    \inheritsavedanchors[from=rectangle]
    \inheritanchorborder[from=rectangle]
    \inheritanchor[from=rectangle]{center}
    \anchor{left}{\pgfpoint{-0.5cm}{0cm}}
    \anchor{right}{\pgfpoint{0.5cm}{0cm}}
    \backgroundpath{
        % Draw the switch shape (rectangle with diagonal lines)
        \pgfsetlinewidth{0.8pt}
        \pgfpathrectanglecorners{\pgfpoint{-0.5cm}{-0.5cm}}{\pgfpoint{0.5cm}{0.5cm}}
        \pgfpathmoveto{\pgfpoint{-0.25cm}{-0.25cm}}
        \pgfpathlineto{\pgfpoint{0cm}{-0.25cm}}
        \pgfpathlineto{\pgfpoint{0cm}{0.25cm}}
        \pgfpathlineto{\pgfpoint{0.25cm}{0.25cm}}
        \pgfusepath{draw}
    }
}
\pgfdeclareshape{quantamp}{
    \inheritsavedanchors[from=op amp]
    \inheritanchorborder[from=op amp]
    \inheritanchor[from=op amp]{center}
    \inheritanchor[from=op amp]{+}
    \inheritanchor[from=op amp]{-}
    \inheritanchor[from=op amp]{out}
    \inheritanchor[from=op amp]{up}
    \inheritanchor[from=op amp]{down}
    \backgroundpath{
        \pgfsetlinewidth{0.8pt}
        \pgfpathmoveto{\pgfpoint{-0.5cm}{0.55cm}}
        \pgfpathlineto{\pgfpoint{0.5cm}{0}}
        \pgfpathlineto{\pgfpoint{-0.5cm}{-0.55cm}}
        \pgfpathclose
        \pgfpathmoveto{\pgfpoint{-0.35cm}{-0.2cm}}
        \pgfpathlineto{\pgfpoint{-0.2cm}{-0.2cm}}
        \pgfpathlineto{\pgfpoint{-0.2cm}{0.2cm}}
        \pgfpathlineto{\pgfpoint{-0.05cm}{0.2cm}}
        \pgfusepath{draw}
    }
}
\pgfdeclareshape{decimation}{
    \inheritsavedanchors[from=rectangle]
    \inheritanchorborder[from=rectangle]
    \inheritanchor[from=rectangle]{center}
    \anchor{left}{\pgfpoint{-0.5cm}{0.0cm}}
    \anchor{right}{\pgfpoint{0.5cm}{0.0cm}}
    \anchor{top}{\pgfpoint{0.0cm}{0.5cm}}
    \anchor{bot}{\pgfpoint{0.0cm}{-0.5cm}}
    \backgroundpath{
        \pgfsetlinewidth{0.8pt}
        \pgfpathrectanglecorners{\pgfpoint{-0.5cm}{-0.5cm}}{\pgfpoint{0.5cm}{0.5cm}}
        \pgfpathmoveto{\pgfpoint{-0.25cm}{0.25cm}}
        \pgfpathlineto{\pgfpoint{-0.25cm}{-0.25cm}}
        \pgfpathmoveto{\pgfpoint{-0.30cm}{-0.185cm}}
        \pgfpathlineto{\pgfpoint{-0.25cm}{-0.25cm}}
        \pgfpathlineto{\pgfpoint{-0.20cm}{-0.185cm}}
        \pgfusepath{draw}
    }
}
\makeatother

\geometry{
    a4paper,
    hmargin=2.54cm,
    tmargin=1.27cm,
    bmargin=1.27cm,
    includeheadfoot
}
\setcounter{secnumdepth}{0}  % Disable section numbering

\begin{filecontents}{linux-graphics-stack.bib}
@misc{navlgs,
  author       = {{Pengutronix}},
  title        = {{ELCE 2022: Navigating the Linux Graphics Stack}},
  year         = {2022},
  howpublished = {\url{https://youtu.be/WaF7mVhnhtE}},
  note         = {Accessed: 2024-12-25}
}
\end{filecontents}
\addbibresource{linux-graphics-stack.bib}

\begin{document}
\section{Linux Graphics Stack}

\cite{navlgs}

\subsection{Some Acronyms}
\begin{verbatim}
ls /sys/class/drm/card0-HDMI-A-1
\end{verbatim}

\begin{verbatim}
ls /sys/kernel/debug/dri/0
\end{verbatim}

\begin{verbatim}
echo 0x1ff > /sys/module/drm/parameters/debug
\end{verbatim}
\begin{itemize}
    \item Wayland
    \item DRI - Direct Rendering Infrastructure
    \item DRM - Direct Rendering Manager
    \item KMS - Kernel Mode Settings
    \item FB - Frame Buffer
    \item EGL
    \item OpenGL
\end{itemize}

\subsection{Topics}

\begin{itemize}
    \item Pixel buffer onto display
    \item Draw pixel into a pixel buffer
    \item Compose multiple pixel buffers
\end{itemize}

\subsection{Display - Acronyms}

\begin{itemize}
    \item DRI - Direct Rendering Infrastructure

        A framework in the X Window System that enables direct access to
        graphics hardware, allowing applications to perform rendering without
        the overhead of inter-process communication.

    \item DRM - Direct Rendering Manager

        A subsystem of the Linux kernel responsible for managing graphics
        processing units (GPUs), providing an API for user-space programs to
        perform operations like configuring display modes and rendering
        graphics.

    \item KMS - Kernel Mode Settings

        A component of the Direct Rendering Manager (DRM) that moves
        responsibility for display mode setting (such as resolution and color
        depth) from user space to the kernel, improving system stability and
        reducing flickering during mode changes.

    \item FB - Frame Buffer

        A portion of RAM containing a bitmap that drives a video display; in
        Linux, the framebuffer device provides a simple interface to graphics
        hardware, allowing applications to write directly to the display
        buffer.
\end{itemize}

\subsection{Display Stack}

\subsection{Kernel Debugging}

\begin{verbatim}
ls /sys/class/drm/card0-HDMI-A-1
\end{verbatim}

\begin{verbatim}
ls /sys/kernel/debug/dri/0
\end{verbatim}

\begin{verbatim}
echo 0x1ff > /sys/module/drm/parameters/debug
\end{verbatim}

\subsection{GPU - Acronyms}

\begin{itemize}
    \item OpenGL - Open Graphics Library

        A cross-platform API for rendering 2D and 3D vector graphics,
        facilitating hardware-accelerated rendering by interfacing with the GPU.

    \item GLSL - OpenGL Shading Language

        A high-level shading language with C-like syntax, enabling developers to
        write shaders for programmable stages of the graphics pipeline, such as
        vertex and fragment shaders

    \item EGL - Native Platform Graphics Interface API

        Provides an interface between OpenGL and the native windowing system
        for managing rendering surfaces.
\end{itemize}

In the Direct Rendering Manager (DRM) framework, only one application can hold
the DRM master status at a time, granting it exclusive control over display
resources.

Traditional toolkits and applications often lack direct support for Kernel Mode
Setting (KMS), making it challenging for multiple applications to manage display
outputs concurrently.

To address this limitation, a window compositor, such as Wayland, is employed.
The compositor acts as the sole DRM master, managing access to display resources
and facilitating the coordination of multiple applications' rendering requests.

This architecture enables a more efficient and secure management of graphical
resources, allowing applications to render content without requiring direct
control over the display hardware.

\newpage
\printbibliography
\end{document}
